% =============================================================================
% COMP0496 — Programação A
% Aula 3: Herança, Polimorfismo e Composição
% Tema: "O Podrão — Food Truck"
% Prof. Dr. André Yoshiaki Kashiwabara — DCOMP/UFS
% =============================================================================

\documentclass[aspectratio=169]{beamer}
\usetheme{UFS}
\usepackage[utf8]{inputenc}
\usepackage[portuguese]{babel}
\usepackage{amsmath,amssymb}
\usepackage{graphicx}
\usepackage{booktabs}
\usepackage{listings}
\usepackage{xcolor}
\usepackage{tikz}
\usetikzlibrary{positioning, arrows.meta, shapes.geometric}

% ---------------------------------------------------------------------------
% lstlisting style — VS Code Light+ palette
% ---------------------------------------------------------------------------
\definecolor{bgcolor}{HTML}{FFFFFF}
\definecolor{linecolor}{HTML}{DDDDDD}
\definecolor{numcolor}{HTML}{999999}
\definecolor{kwcolor}{HTML}{0000FF}
\definecolor{strcolor}{HTML}{A31515}
\definecolor{cmtcolor}{HTML}{008000}
\definecolor{fncolor}{HTML}{795E26}
\definecolor{typecolor}{HTML}{267F99}

\lstdefinestyle{modernstyle}{
  language=Python,
  backgroundcolor=\color{bgcolor},
  rulecolor=\color{linecolor},
  numberstyle=\tiny\color{numcolor},
  keywordstyle=\color{kwcolor}\bfseries,
  stringstyle=\color{strcolor},
  commentstyle=\color{cmtcolor}\itshape,
  emph={print,len,range,super,isinstance,type,abs},
  emphstyle=\color{fncolor},
  emph={[2]int,float,str,list,dict,bool,None,True,False},
  emphstyle={[2]\color{typecolor}},
  basicstyle=\ttfamily\scriptsize,
  frame=single,
  numbers=left,
  numbersep=8pt,
  tabsize=4,
  showstringspaces=false,
  breaklines=true,
  columns=fullflexible,
  extendedchars=true,
  literate={á}{{\'a}}1 {é}{{\'e}}1 {í}{{\'i}}1 {ó}{{\'o}}1 {ú}{{\'u}}1
           {Á}{{\'A}}1 {É}{{\'E}}1 {Í}{{\'I}}1 {Ó}{{\'O}}1 {Ú}{{\'U}}1
           {à}{{\`a}}1 {è}{{\`e}}1 {ì}{{\`i}}1 {ò}{{\`o}}1 {ù}{{\`u}}1
           {ã}{{\~a}}1 {õ}{{\~o}}1 {Ã}{{\~A}}1 {Õ}{{\~O}}1
           {ç}{{\c{c}}}1 {Ç}{{\c{C}}}1
           {ê}{{\^e}}1 {â}{{\^a}}1 {ô}{{\^o}}1 {î}{{\^i}}1 {û}{{\^u}}1
           {Ê}{{\^E}}1 {Â}{{\^A}}1 {Ô}{{\^O}}1
           {—}{{---}}1 {–}{{--}}1 {…}{{\ldots}}1
           {“}{{``}}1 {”}{{''}}1 {‘}{{`}}1 {’}{{'}}1,
}
\lstset{style=modernstyle}

% ---------------------------------------------------------------------------
% Metadata
% ---------------------------------------------------------------------------
\title{COMP0496 -- Programação A \\ Programação Orientada a Objetos — Parte 3}
\subtitle{Herança, Polimorfismo e Composição}
\author{Prof.\ Dr.\ André Yoshiaki Kashiwabara}
\institute{DCOMP -- Universidade Federal de Sergipe\\\texttt{andreyoshiaki@dcomp.ufs.br}}
\date{}

\begin{document}

% ===== TÍTULO =====
\begin{frame}[plain,shrink]
  \titlepage
\end{frame}

% ===== ROTEIRO =====
\begin{frame}{Roteiro}
  \tableofcontents
\end{frame}

% =============================================================================
% SECTION 1: O Problema — Código Duplicado
% =============================================================================
\section{O Problema — Código Duplicado}

\begin{frame}[fragile]{Três Classes Independentes}
  No food truck \textbf{O Podrão}, temos três lanches:

  \begin{columns}[T]
    \column{0.32\textwidth}
    \begin{lstlisting}
class XTudo:
    def __init__(self, preco):
        self.preco = preco
        self.nome = "X-Tudo"

    def descrever(self):
        return f"{self.nome}: R${self.preco:.2f}"

    def preparar(self):
        print("Fritando hambúrguer...")
    \end{lstlisting}

    \column{0.32\textwidth}
    \begin{lstlisting}
class Dogao:
    def __init__(self, preco):
        self.preco = preco
        self.nome = "Dogão"

    def descrever(self):
        return f"{self.nome}: R${self.preco:.2f}"

    def preparar(self):
        print("Montando hot-dog...")
    \end{lstlisting}

    \column{0.32\textwidth}
    \begin{lstlisting}
class Acai:
    def __init__(self, preco):
        self.preco = preco
        self.nome = "Açaí"

    def descrever(self):
        return f"{self.nome}: R${self.preco:.2f}"

    def preparar(self):
        print("Batendo açaí...")
    \end{lstlisting}
  \end{columns}
\end{frame}

\begin{frame}{Por que Isso é Ruim?}
  \begin{alertblock}{Violação do Princípio DRY}
    \textbf{Don't Repeat Yourself} — código duplicado é fonte de bugs.
  \end{alertblock}

  \bigskip

  \begin{itemize}
    \item \texttt{\_\_init\_\_} e \texttt{descrever} são \textbf{idênticos} nas três classes
    \item Se precisar adicionar \texttt{calorias}: alterar \textbf{três} classes
    \item Se corrigir um bug em \texttt{descrever}: corrigir em \textbf{três} lugares
    \item E se o cardápio crescer para 15 itens?
  \end{itemize}

  \bigskip

  \begin{exampleblock}{Solução}
    Extrair o que é \textbf{comum} para uma classe base $\Rightarrow$ \textbf{Herança}.
  \end{exampleblock}
\end{frame}

% =============================================================================
% SECTION 2: Herança Simples
% =============================================================================
\section{Herança Simples}

\begin{frame}{O Relacionamento ``É Um''}
  Herança modela o relacionamento \textbf{``é um''} (\emph{is-a}):

  \bigskip

  \begin{center}
    \begin{tabular}{lcc}
      \toprule
      \textbf{Afirmação} & \textbf{Verdade?} & \textbf{Herança?} \\
      \midrule
      XTudo é um Lanche           & \checkmark & Sim \\
      Dogão é um Lanche           & \checkmark & Sim \\
      Açaí é um Lanche            & \checkmark & Sim \\
      Pedido é uma CaixaRegistradora & $\times$ & Não \\
      Pedido \emph{tem} Lanches   & \checkmark & Composição \\
      \bottomrule
    \end{tabular}
  \end{center}

  \bigskip

  \begin{alertblock}{Regra de Ouro}
    Se ``B é um A'' faz sentido no domínio $\Rightarrow$ \texttt{class B(A)}.
  \end{alertblock}
\end{frame}

\begin{frame}{Hierarquia de Classes}
  % Coordinate map:
  % (0,3)    = Lanche (parent node)
  % (-3,0)   = XTudo (child left)
  % (0,0)    = Dogao (child center)
  % (3,0)    = Acai (child right)
  \begin{center}
    \begin{tikzpicture}[scale=1.1, every node/.style={scale=1.1}]
      \node[draw, fill=blue!15, rounded corners, minimum width=3cm, minimum height=1cm, font=\bfseries]
        (lanche) at (0,3) {Lanche};
      \node[draw, fill=green!10, rounded corners, minimum width=2.5cm, minimum height=1cm]
        (xtudo) at (-4,0) {XTudo};
      \node[draw, fill=green!10, rounded corners, minimum width=2.5cm, minimum height=1cm]
        (dogao) at (0,0) {Dogão};
      \node[draw, fill=green!10, rounded corners, minimum width=2.5cm, minimum height=1cm]
        (acai) at (4,0) {Açaí};

      \draw[->, thick] (xtudo) -- node[left, font=\scriptsize]{herda} (lanche);
      \draw[->, thick] (dogao) -- node[right, font=\scriptsize]{herda} (lanche);
      \draw[->, thick] (acai) -- node[right, font=\scriptsize]{herda} (lanche);

      \node[below=0.2cm of xtudo, font=\scriptsize\ttfamily, text=gray] {class XTudo(Lanche)};
      \node[below=0.2cm of dogao, font=\scriptsize\ttfamily, text=gray] {class Dogao(Lanche)};
      \node[below=0.2cm of acai, font=\scriptsize\ttfamily, text=gray] {class Acai(Lanche)};
    \end{tikzpicture}
  \end{center}
\end{frame}

\begin{frame}[fragile]{Classe Base: \texttt{Lanche}}
  \begin{lstlisting}
class Lanche:
    def __init__(self, nome, preco):
        self.nome = nome
        self.preco = preco

    def descrever(self):
        return f"{self.nome}: R${self.preco:.2f}"

    def preparar(self):
        print(f"Preparando {self.nome}...")
  \end{lstlisting}

  \begin{exampleblock}{O que mudou?}
    Todo código comum agora vive em \textbf{um único lugar}.
  \end{exampleblock}
\end{frame}

\begin{frame}[fragile]{Subclasses com \texttt{super()}}
  \begin{lstlisting}
class XTudo(Lanche):
    def __init__(self, preco, molho="especial"):
        super().__init__("X-Tudo", preco)
        self.molho = molho

    def preparar(self):
        print(f"Fritando hambúrguer com molho {self.molho}...")
  \end{lstlisting}

  \begin{lstlisting}
class Dogao(Lanche):
    def __init__(self, preco, tamanho="grande"):
        super().__init__("Dogão", preco)
        self.tamanho = tamanho

    def preparar(self):
        print(f"Montando hot-dog {self.tamanho}...")
  \end{lstlisting}
\end{frame}

\begin{frame}[fragile]{Herdando Métodos Automaticamente}
  \begin{lstlisting}
x = XTudo(25.0)
d = Dogao(18.0)

# descrever() é herdado de Lanche — funciona sem reescrever!
print(x.descrever())   # X-Tudo: R$25.00
print(d.descrever())   # Dogão: R$18.00

# preparar() foi sobrescrito em cada subclasse
x.preparar()   # Fritando hambúrguer com molho especial...
d.preparar()   # Montando hot-dog grande...
  \end{lstlisting}

  \begin{exampleblock}{Regra}
    Se a subclasse \textbf{não} redefine um método, usa o da superclasse.
  \end{exampleblock}
\end{frame}

\begin{frame}[fragile]{Verificando Tipos com \texttt{isinstance()}}
  \begin{lstlisting}
x = XTudo(25.0)

isinstance(x, XTudo)    # True
isinstance(x, Lanche)   # True  (XTudo É UM Lanche)
isinstance(x, Dogao)    # False

type(x) == XTudo         # True
type(x) == Lanche        # False (type() é exato)
  \end{lstlisting}

  \bigskip

  \begin{center}
    \begin{tabular}{lcc}
      \toprule
      \textbf{Verificação} & \textbf{\texttt{isinstance}} & \textbf{\texttt{type() ==}} \\
      \midrule
      Considera herança? & Sim & Não \\
      Uso recomendado    & Geral & Raro \\
      \bottomrule
    \end{tabular}
  \end{center}
\end{frame}

% =============================================================================
% SECTION 3: Sobrescrita de Métodos
% =============================================================================
\section{Sobrescrita de Métodos}

\begin{frame}[fragile]{Override: Redefinindo um Método}
  \textbf{Sobrescrita} (\emph{override}): a subclasse fornece sua própria versão de um método herdado.

  \begin{lstlisting}
class Acai(Lanche):
    def __init__(self, preco, complementos=None):
        super().__init__("Açaí", preco)
        self.complementos = complementos or ["granola", "banana"]

    def preparar(self):
        # Sobrescreve completamente o método de Lanche
        print(f"Batendo açaí com {', '.join(self.complementos)}...")
  \end{lstlisting}

  \begin{alertblock}{Atenção}
    O método da subclasse \textbf{substitui} o da superclasse — não são acumulados.
  \end{alertblock}
\end{frame}

\begin{frame}[fragile]{Padrão Extend: \texttt{super()} + Extensão}
  Quando queremos \textbf{estender} (não substituir) o comportamento do pai:

  \begin{lstlisting}
class XTudoEspecial(XTudo):
    def __init__(self, preco, extras=None):
        super().__init__(preco, molho="da casa")
        self.extras = extras or ["bacon", "ovo"]

    def descrever(self):
        base = super().descrever()  # chama Lanche.descrever()
        return f"{base} + {', '.join(self.extras)}"
  \end{lstlisting}

  \begin{lstlisting}
xe = XTudoEspecial(32.0)
print(xe.descrever())
# X-Tudo: R$32.00 + bacon, ovo
  \end{lstlisting}
\end{frame}

\begin{frame}[fragile]{Armadilha: Esquecer \texttt{super().\_\_init\_\_()}}
  \begin{lstlisting}
class Coxinha(Lanche):
    def __init__(self, preco):
        # ESQUECEU super().__init__() !!!
        self.recheio = "frango"

c = Coxinha(8.0)
c.recheio       # "frango" — OK
c.descrever()   # AttributeError: 'Coxinha' has no attribute 'nome'
  \end{lstlisting}

  \begin{alertblock}{Por quê?}
    Sem \texttt{super().\_\_init\_\_()}, os atributos \texttt{nome} e \texttt{preco} nunca são criados.
  \end{alertblock}

  \begin{exampleblock}{Correção}
    Sempre chame \texttt{super().\_\_init\_\_()} no início do \texttt{\_\_init\_\_} da subclasse.
  \end{exampleblock}
\end{frame}

% =============================================================================
% SECTION 4: Polimorfismo e Duck Typing
% =============================================================================
\section{Polimorfismo e Duck Typing}

\begin{frame}[fragile]{Polimorfismo: Mesma Chamada, Comportamentos Diferentes}
  \textbf{Polimorfismo}: objetos de tipos diferentes respondem à mesma mensagem, cada um à sua maneira.

  \begin{lstlisting}
pedido = [XTudo(25.0), Dogao(18.0), Acai(15.0)]

for item in pedido:
    item.preparar()  # cada classe executa seu próprio preparar()
  \end{lstlisting}

  Saída:
  \begin{lstlisting}[numbers=none]
Fritando hambúrguer com molho especial...
Montando hot-dog grande...
Batendo açaí com granola, banana...
  \end{lstlisting}

  \begin{exampleblock}{Essência}
    O código \textbf{não precisa saber} o tipo exato — basta que o objeto tenha o método.
  \end{exampleblock}
\end{frame}

\begin{frame}[fragile]{Duck Typing}
  \begin{quote}
    \emph{``If it walks like a duck and quacks like a duck, then it must be a duck.''}
  \end{quote}

  \bigskip

  \begin{lstlisting}
class Bebida:          # NÃO herda de Lanche!
    def __init__(self, nome, preco):
        self.nome = nome
        self.preco = preco

    def preparar(self):
        print(f"Servindo {self.nome}...")

    def descrever(self):
        return f"{self.nome}: R${self.preco:.2f}"
  \end{lstlisting}
\end{frame}

\begin{frame}[fragile,shrink]{Duck Typing na Prática}
  \begin{lstlisting}
pedido = [XTudo(25.0), Dogao(18.0), Bebida("Guaraná", 6.0)]

for item in pedido:
    print(item.descrever())
    item.preparar()
  \end{lstlisting}

  Saída:
  \begin{lstlisting}[numbers=none]
X-Tudo: R$25.00
Fritando hambúrguer com molho especial...
Dogão: R$18.00
Montando hot-dog grande...
Guaraná: R$6.00
Servindo Guaraná...
  \end{lstlisting}

  \begin{alertblock}{Observação}
    \texttt{Bebida} não herda de \texttt{Lanche}, mas funciona no mesmo loop porque tem os mesmos métodos. Isso é \textbf{duck typing}.
  \end{alertblock}
\end{frame}

% =============================================================================
% SECTION 5: Composição — A Alternativa
% =============================================================================
\section{Composição — A Alternativa}

\begin{frame}{Herança vs.\ Composição}
  \begin{center}
    \begin{tabular}{lll}
      \toprule
      \textbf{Relacionamento} & \textbf{Mecanismo} & \textbf{Exemplo} \\
      \midrule
      ``É um'' (\emph{is-a})  & Herança    & XTudo \textbf{é um} Lanche \\
      ``Tem um'' (\emph{has-a}) & Composição & Pedido \textbf{tem} Lanches \\
      ``Tem um'' (\emph{has-a}) & Composição & Pedido \textbf{tem} FormaPagamento \\
      ``É um'' (\emph{is-a})  & Herança    & Pix \textbf{é uma} FormaPagamento \\
      \bottomrule
    \end{tabular}
  \end{center}
\end{frame}

\begin{frame}[fragile,shrink]{Composição no Código}
  \begin{lstlisting}
class FormaPagamento:
    def processar(self, valor):
        raise NotImplementedError

class Pix(FormaPagamento):
    def processar(self, valor):
        print(f"Pix de R${valor:.2f} confirmado!")

class Pedido:
    def __init__(self, pagamento):
        self.itens = []               # composição: tem lanches
        self.pagamento = pagamento    # composição: tem forma de pgto

    def adicionar(self, item):
        self.itens.append(item)

    def fechar(self):
        total = sum(i.preco for i in self.itens)
        self.pagamento.processar(total)
  \end{lstlisting}
\end{frame}

\begin{frame}[fragile]{Usando o Pedido}
  \begin{lstlisting}
pix = Pix()
pedido = Pedido(pix)
pedido.adicionar(XTudo(25.0))
pedido.adicionar(Bebida("Guaraná", 6.0))
pedido.fechar()
# Pix de R$31.00 confirmado!
  \end{lstlisting}

  \begin{exampleblock}{Gang of Four}
    \emph{``Favor composition over inheritance.''} \\
    Se você só quer \textbf{reusar} comportamento, sem relação conceitual ``é um'', use composição.
  \end{exampleblock}
\end{frame}

\begin{frame}{Quando Usar Cada Um?}
  \begin{center}
    \begin{tabular}{lcc}
      \toprule
      \textbf{Critério} & \textbf{Herança} & \textbf{Composição} \\
      \midrule
      Relação ``é um'' genuína & \checkmark & \\
      Apenas reusar código     &   & \checkmark \\
      Trocar comportamento em runtime & & \checkmark \\
      Hierarquia natural e estável & \checkmark & \\
      Flexibilidade máxima & & \checkmark \\
      \bottomrule
    \end{tabular}
  \end{center}

  \begin{alertblock}{Regra Prática}
    Comece com composição. Use herança apenas quando o ``é um'' for claro e estável.
  \end{alertblock}
\end{frame}

% =============================================================================
% SECTION 6: Herança Múltipla e MRO
% =============================================================================
\section{Herança Múltipla e MRO}

\begin{frame}[fragile]{Herança Múltipla em Python}
  Python permite herdar de \textbf{mais de uma classe}:

  \begin{lstlisting}
class Promocional:
    def aplicar_desconto(self, pct):
        self.preco *= (1 - pct / 100)

class XTudoPromo(XTudo, Promocional):
    pass

xp = XTudoPromo(25.0)
xp.aplicar_desconto(10)
print(xp.preco)   # 22.5
  \end{lstlisting}
\end{frame}

\begin{frame}[fragile]{MRO — Method Resolution Order}
  Quando há conflito, Python usa o \textbf{C3 Linearization} para decidir qual método chamar:

  \begin{lstlisting}
print(XTudoPromo.__mro__)
# (XTudoPromo, XTudo, Lanche, Promocional, object)
  \end{lstlisting}

  Python busca o método na ordem da MRO: da esquerda para a direita.

  \begin{alertblock}{Conselho}
    Herança múltipla pode causar confusão. Na prática:
    \begin{itemize}
      \item Use \textbf{mixins} simples (como \texttt{Promocional})
      \item Prefira \textbf{composição} para combinações complexas
      \item Evite hierarquias diamante
    \end{itemize}
  \end{alertblock}
\end{frame}

% =============================================================================
% SECTION 7: Classes Abstratas
% =============================================================================
\section{Classes Abstratas}

\begin{frame}[fragile,shrink]{O Problema: Métodos ``Esquecidos''}
  \begin{lstlisting}
class Funcionario:
    def __init__(self, nome, salario_base):
        self.nome = nome
        self.salario_base = salario_base

    def calcular_salario(self):
        raise NotImplementedError("Subclasse deve implementar")
  \end{lstlisting}

  \begin{lstlisting}
class Atendente(Funcionario):
    pass   # esqueceu de implementar calcular_salario!

a = Atendente("Maria", 1500)
a.calcular_salario()  # NotImplementedError — mas só descobre AQUI!
  \end{lstlisting}

  \begin{alertblock}{Problema}
    O erro só aparece quando \texttt{calcular\_salario()} é chamado — pode ser tarde demais.
  \end{alertblock}
\end{frame}

\begin{frame}[fragile,shrink]{Solução: \texttt{ABC} e \texttt{@abstractmethod}}
  \begin{lstlisting}
from abc import ABC, abstractmethod

class Funcionario(ABC):
    def __init__(self, nome, salario_base):
        self.nome = nome
        self.salario_base = salario_base

    @abstractmethod
    def calcular_salario(self):
        """Cada tipo de funcionário calcula diferente."""
        pass
  \end{lstlisting}

  \begin{lstlisting}
# Tentar instanciar a classe abstrata:
f = Funcionario("João", 2000)
# TypeError: Can't instantiate abstract class Funcionario
#   with abstract method calcular_salario
  \end{lstlisting}

  \begin{exampleblock}{Vantagem}
    O erro ocorre na \textbf{instanciação}, não na chamada do método — falha cedo!
  \end{exampleblock}
\end{frame}

\begin{frame}[fragile,shrink]{Subclasses Concretas}
  \begin{lstlisting}
class Atendente(Funcionario):
    def calcular_salario(self):
        return self.salario_base

class Cozinheiro(Funcionario):
    def __init__(self, nome, salario_base, adicional_insalubridade):
        super().__init__(nome, salario_base)
        self.adicional = adicional_insalubridade

    def calcular_salario(self):
        return self.salario_base + self.adicional

class Entregador(Funcionario):
    def __init__(self, nome, salario_base, entregas):
        super().__init__(nome, salario_base)
        self.entregas = entregas

    def calcular_salario(self):
        return self.salario_base + self.entregas * 5
  \end{lstlisting}
\end{frame}

\begin{frame}{Por que Classes Abstratas Importam}
  \begin{center}
    \small
    \begin{tabular}{lcc}
      \toprule
      \textbf{Aspecto} & \texttt{raise NotImplementedError} & \texttt{ABC + @abstractmethod} \\
      \midrule
      Quando dá erro   & Na chamada do método & Na instanciação \\
      Documentação      & Implícita             & Explícita \\
      Forçar implementação & Não & Sim \\
      \bottomrule
    \end{tabular}
  \end{center}

  \begin{exampleblock}{Princípio}
    Classes abstratas definem um \textbf{contrato}: ``toda subclasse deve ter estes métodos''.
  \end{exampleblock}
\end{frame}

% =============================================================================
% SECTION 8: Prática — Folha de Pagamento
% =============================================================================
\section{Prática — Folha de Pagamento}

\begin{frame}{Cenário: Folha de Pagamento do Food Truck}
  O Podrão precisa calcular a folha de pagamento mensal.

  \bigskip

  \begin{center}
    \begin{tabular}{lll}
      \toprule
      \textbf{Cargo} & \textbf{Cálculo} & \textbf{Exemplo} \\
      \midrule
      Atendente   & Salário base              & R\$\,1.500 \\
      Cozinheiro  & Base + insalubridade      & R\$\,1.800 + R\$\,300 \\
      Entregador  & Base + R\$\,5/entrega     & R\$\,1.200 + 80 $\times$ R\$\,5 \\
      \bottomrule
    \end{tabular}
  \end{center}
\end{frame}

\begin{frame}[fragile,shrink]{Loop Polimórfico para a Folha}
  \begin{lstlisting}
equipe = [
    Atendente("Maria", 1500),
    Cozinheiro("João", 1800, 300),
    Entregador("Carlos", 1200, 80),
    Atendente("Ana", 1500),
]

total_folha = 0
for func in equipe:
    salario = func.calcular_salario()
    total_folha += salario
    print(f"{func.nome:.<20} R${salario:>8.2f}")

print(f"{'TOTAL':.<20} R${total_folha:>8.2f}")
  \end{lstlisting}

  Saída:
  \begin{lstlisting}[numbers=none]
Maria............... R$ 1500.00
João................ R$ 2100.00
Carlos.............. R$ 1600.00
Ana................. R$ 1500.00
TOTAL............... R$ 6700.00
  \end{lstlisting}
\end{frame}

\begin{frame}[fragile]{O Poder do Polimorfismo}
  Observe o loop:

  \begin{lstlisting}
for func in equipe:
    salario = func.calcular_salario()
  \end{lstlisting}

  \bigskip

  \begin{itemize}
    \item O loop \textbf{não sabe} se \texttt{func} é Atendente, Cozinheiro ou Entregador
    \item Cada objeto ``sabe'' calcular seu próprio salário
    \item Adicionar novo cargo (ex: \texttt{Gerente}) não altera o loop
    \item Isso é o \textbf{Open/Closed Principle}: aberto para extensão, fechado para modificação
  \end{itemize}

  \begin{exampleblock}{Reflexão}
    Polimorfismo + classes abstratas = código extensível e seguro.
  \end{exampleblock}
\end{frame}

% =============================================================================
% SECTION 9: Resumo
% =============================================================================
\section{Resumo}

\begin{frame}{Tabela Resumo}
  \begin{center}
    \small
    \begin{tabular}{llp{6cm}}
      \toprule
      \textbf{Conceito} & \textbf{Palavra-chave} & \textbf{Quando usar} \\
      \midrule
      Herança simples   & \texttt{class B(A)} & Relação ``é um'' clara \\
      \texttt{super()}  & \texttt{super().\_\_init\_\_()} & Delegar ao construtor pai \\
      Sobrescrita       & Override             & Especializar comportamento \\
      Polimorfismo      & Mesma interface      & Tratar objetos uniformemente \\
      Duck typing       & (implícito)          & Flexibilidade sem herança \\
      Composição        & Atributo de instância & Relação ``tem um'' \\
      Herança múltipla  & \texttt{class C(A,B)} & Mixins simples \\
      Classes abstratas & \texttt{ABC}, \texttt{@abstractmethod} & Forçar contrato \\
      \bottomrule
    \end{tabular}
  \end{center}
\end{frame}

\begin{frame}{Princípios-Chave}
  \begin{enumerate}
    \item \textbf{DRY}: não duplique código — extraia para a superclasse
    \item \textbf{Liskov}: subclasse deve ser usável onde a superclasse é esperada
    \item \textbf{Composição sobre herança}: use herança apenas para ``é um'' genuíno
    \item \textbf{Falhe cedo}: prefira \texttt{ABC} a \texttt{NotImplementedError}
    \item \textbf{Polimorfismo}: programe para a interface, não para a implementação
  \end{enumerate}

  \bigskip

  \begin{exampleblock}{Próxima Aula (Aula 4)}
    \texttt{@dataclass}, princípios SOLID e introdução a Design Patterns.
  \end{exampleblock}
\end{frame}

\end{document}
